\section{Main results on the frozen test set}
\label{sec:results}

\generated{table_main}

\generated{table_main_speedup}

\generated{table_agreement}

The comparison is between runs that produce the same object: a certified primal
solution, the multipliers, and the canonical blocks of the positive path at
the instance's own capacity. Instances that no configuration solved within the
limit are still counted in PAR-2 and in the solved column; they are not removed
from the sample.

\begin{figure}[tbp]
  \centering
  \IfFileExists{generated/profile_main.dat}{%
    \begin{tikzpicture}
      \begin{axis}[xlabel={$\tau$ (factor over the best time on the instance)},
                   ylabel={fraction of instances}, xmode=log, ymin=0, ymax=1.02,
                   legend pos=south east, width=0.62\textwidth, height=0.42\textwidth]
        \addplot table[x=tau, y index=1] {generated/profile_main.dat};
        \addplot table[x=tau, y index=2] {generated/profile_main.dat};
        \addplot table[x=tau, y index=3] {generated/profile_main.dat};
        \addplot table[x=tau, y index=4] {generated/profile_main.dat};
        \legend{\texttt{dinkelbach},\texttt{forest-reference},\texttt{forest-tuned},\texttt{parametric}}
      \end{axis}
    \end{tikzpicture}}{\small\itshape Performance profile not available: the main
    campaign has not been analysed in this build of the report.}
  \caption{Performance profile on the frozen test set. A configuration is
  counted at $\tau$ when its median time on an instance is within a factor
  $\tau$ of the best median time recorded on that instance; unsolved instances
  never enter the numerator. The compared set is the instances at least one
  configuration solved and verified.}
  \label{fig:profile-main}
\end{figure}

\subsection{Against the 2023 prototype}
\label{sec:legacy}

The group's 2023 \FS{} prototype was imported into this project under
\file{legacy/forest2023/} and measured on the same maximum-ratio oracle, which is
the task it implements. Three patches were needed to run it at all and are
recorded in that directory: an absolute CPLEX include replaced by opaque
typedefs, four interactive \texttt{cin.get()} pauses removed, and the append to a
file in the working directory removed. No line of the algorithm was changed.

Running it exposed limits that belong to the prototype, not to the comparison.
It stores profits and weights as \texttt{int}, so instances with fractional data
cannot be represented and are refused instead of rounded. It allocates two
$n\times n$ adjacency matrices, so its memory is quadratic in the number of
vertices. Its initial basis assumes a spanning tree of $n-1$ edges, so it requires a
connected graph. Finally, the feasibility shortcut added in the 2023-05-11
version recurses without termination on a generic, non-bipartite DAG --- it
overflows the stack already on the eight-node example --- so the measurements use
it disabled, which reproduces the behaviour of the 2023-05-05 version of the same
file.

\paragraph{The outcome.} On the eighteen instances measured, the prototype
completed \emph{none}, so there is no time to compare it on and no table worth
printing. Four times it hit the time limit --- including a chain of two hundred
vertices, which the current solver certifies in nine milliseconds --- once it
crashed, once it exceeded the memory budget, and the remaining twelve it refused
outright, either because the data are not integral or because the graph is
disconnected and no spanning tree of $n-1$ edges exists. The current solver
returned a certified optimal ratio on every one of the eighteen, in times between
seven milliseconds and thirty-two seconds. The per-instance record is in
\file{experiments/analysis/legacy\_baseline*.json}.

That is a wide gap, and it should be read for what it is. Most of it is not a
difference in pivot rules: it is that the prototype rebuilds the whole forest and
reallocates two arrays at every iteration, recurses instead of iterating, and
carries data structures quadratic in the number of vertices. Those are the costs of
a research prototype that was never asked to run outside the bipartite instances
it was written for, and the current solver inherits none of them. The comparison
therefore measures engineering, not the merit of the underlying method --- which
is the same method in both.

What the prototype did contribute is two ideas that survived measurement and are
in the current solver: the early exit of the ratio test on a zero-step blocker,
which it implemented as an explicit feasibility check during the forest
traversal, and the absence of a permanently engaged anti-cycling rule, which the
ablation of section~\ref{sec:analysis} confirms was costing a factor in pivots. A
third idea is still on the table and not yet implemented: the Euler-tour
intervals it maintained, which give subtree membership in constant time and would
attack the rebuild cost that now dominates.

Where both solvers return an optimal ratio, they return the same one. On this
instance set that condition never held, so the cross-check rests on the
eight-node example of the tutorial, where both return the ratio $2$.

